% ADDRESSED - Updated figure 1
\feedback{This [PythonTutor] figure is unreadable.}
{}

% ADDRESSED - Spelling corrected
\feedback{[PythonTutor] Spelling.}
{}

% ADDRESSED - Added supported languages
\feedback{Which languages [are supported by PythonTutor]?}
{}

% ADDRESSED - Added citation for PythonTutor
\feedback{You should have a citation for where to access the [PythonTutor] tool itself and not just the paper describing its efficacy.}
{}

% ADDRESSED - Updated Figure 2
\feedback{The text in the [control-flow graph] figure is too small.}
{}

% ADDRESSED - Added citation for Graphviz
\feedback{[Graphviz] Linky please.}
{}

% ADDRESSED - Added citation for Manim
\feedback{Linky to Manim.}
{}

% ADDRESSED - Removed (he didn't win a Streamy)
\feedback{citation needed [for 3Blue1Brown's winning a Streamy award]}
{}

% ADDRESSED - Updated order of subsections
\feedback{The order of the [model] diagram is in the opposite order that they are presented in [implementation plan] text.}
{}

% ADDRESSED - Increased timeline detail
\feedback{The items in the timeline are incredibly brief and need more detail or support.}
{}

\feedbackdivider

% ADDRESSED - Updated figures
\feedback{The figures are quite shrunk down especially figure 1 and figure 2 making it extremely difficult to read the content within it, if those were made larger it would greatly improve the readability.}
{}

% REFUTED
\feedback{The abstract could briefly mention the intended users explicitly, such as educators, students, or software engineers. Additionally, the abstract could briefly summarize how LLVManim achieves this goal rather than only stating the objective.}
{
  The intended users of LLVManim is a broad spectrum, including educators, students, and software engineers. We believe the abstract already addresses the context in which LLVManim is useful, and to specify the users that exist within those contexts would be redundant.
}

% ADDRESSED - Updated introduction
\feedback{The introduction could more explicitly define the scope of the LLVManim system itself earlier in the section. Currently, LLVManim is only briefly mentioned at the end. Introducing LLVManim sooner would improve clarity and help frame the remainder of the document.}
{}

% REFUTED
\feedback{[The PythonTutor] section could be slightly condensed. Some explanations regarding educational benefits are repeated, and the limitations relevant specifically to LLVManim's goals like their lack of customization, presentation-quality output, etc., should be emphasized more clearly. Additionally, explicitly comparing PythonTutor directly to LLVManim in one or two sentences would make the relevance stronger.}
{
  We believe the current level of detail in the PythonTutor section is appropriate for providing necessary context and motivation for LLVManim. We may provide a more explicit comparison to LLVManim in the future, but we feel the current description sufficiently highlights the limitations of PythonTutor that LLVManim aims to address.
}

% REFUTED
\feedback{[The LLVM Analysis and Transform Passes] section could be improved by separating the strengths and weaknesses into clearly defined subsections or paragraphs. Currently, strengths and limitations are interwoven, which makes it slightly harder to skim. Additionally, some technical explanations like the IR syntax examples could be shortened slightly, as the focus of the proposal is visualization rather than compiler internals.}
{
  We believe the current structure of the LLVM Analysis and Transform Passes section effectively conveys the necessary information without overcomplicating the overall document structure. It follows the same format as the other sections describing existing tools, and we feel it is consistent with the overall style of the proposal.
}

% ADDRESSED/REFUTED - Reduced historical background
\feedback{[The Manim] section could more explicitly explain how Manim integrates into the LLVManim system pipeline. Additionally, some historical background about 3Blue1Brown could be slightly condensed to maintain focus on the project itself.}
{
  We reduced the historical background about 3Blue1Brown. However, these sections are focused on providing context for the tools we are using, and we feel that the current level of detail is appropriate.
}

% ADDRESSED - Removed trace mapping and usability metrics from goals, combined core and major features
\feedback
{Some proof-of-concept metrics (such as $\geq95\%$ trace mapping and $\geq80\%$ usability success rate) could briefly explain how these metrics will be measured or evaluated. This would strengthen the feasibility and evaluation clarity. Additionally, the distinction between core features and major features could be clarified slightly, as both appear required for the full system.}
{}

% ADDRESSED - Aligned Implementation Plan points with architecture diagram
\feedback{The architecture diagram should be explicitly referenced and described more clearly in the text. Explaining what each arrow and component represents would improve clarity. Additionally, some implementation details (such as exact timing offsets and layout algorithms) may be overly detailed for a proposal and could be slightly condensed.}
{}

% ADDRESSED - Improved User Interface section
\feedback{The output format could be explained more clearly, particularly how it connects to the Transformation layer.}
{}


% ADDRESSED - Removed mention of Sugiyama algorithm
\feedback{The scene graph structure could benefit from a brief visual example or mock representation to clarify how data flows through the system. Additionally, the modified Sugiyama layout algorithm is mentioned but not briefly explained, which may confuse readers unfamiliar with graph layout techniques.}
{}

% ADDRESSED - Improved Transformation/Presentation Layer sections
\feedback{The relationship between scene graph generation and Manim output could be explained more clearly in terms of data flow. Additionally, it could briefly explain expected output formats and user workflows.}
{}

% ADDRESSED - Improved error handling explanation
\feedback{[The User Interface] section could briefly explain error handling and limitations more clearly, such as unsupported language features or missing debug metadata. Additionally, this section could clarify whether the CLI is intended for developers, educators, or both.}
{}

% REFUTED
\feedback{Testing and validation phases could be more clearly defined, particularly usability testing and feature validation. Additionally, stretch feature implementation is only briefly mentioned and could be explicitly identified in the timeline.}
{
  We have included testing and validation as part of the overall implementation plan, but we have not explicitly defined separate phases for these activities. We will be continuously testing and validating our implementation throughout the development process, and we will prioritize core features before stretch features to ensure we have a MVP by the end of the semester.
}

% REFUTED
\feedback{The feedback section [...] could be improved by explicitly explaining how each feedback item will be addressed moving forward rather than only stating that it remains unaddressed. Additionally, creating mockups or example animation output would significantly strengthen the proposal and demonstrate feasibility.}
{
  We have addressed feedback in the manner prescribed by the project guidelines, which is to either address it or explain why it remains unaddressed.
}

% REFUTED
\feedback{Ensure consistent formatting across all references and verify citation formatting matches IEEE or required course format exactly.}
{
  We have reviewed our references and believe they are formatted correctly according to the IEEE style. If there are any specific formatting issues, we would appreciate further clarification so we can address them appropriately.
}

\feedbackdivider

% REFUTED
\feedback{It's unnecessary to explain how Julia works. It distracts from the main focus of the paper.}
{
  We do not explain how Julia works in the proposal. We have reworded the relevant section for clarity.
}

% ADDRESSED - Removed mention of streamy award
\feedback{You don't have to include all these accomplishments. We already get the idea that [3Blue1Brown] is popular and well respected from earlier.}
{}

% ADDRESSED - Added Manim example (figure 3)
\feedback{Including pictures of Manim diagrams here would illustrate your point better.}
{}

% REFUTED
\feedback{Your goals seem overly ambitious for just a one semester student project.}
{
  We acknowledge that our goals are ambitious, but we believe they are achievable with careful planning and prioritization. We also believe that our approach to the project is more relevant to the goals of the class than the final output, and we are committed to making significant progress towards our objectives regardless of the final scope of the project.
}

% REFUTED
\feedback{It's unrealistic to expect work to be done over spring break.}
{
  We appreciate the concern, but our teammembers are aware of the timeline and have planned accordingly.
}

% ADDRESSED - Added risk evaluation section
\feedback{You should talk more about the potential risks. It's not really mentioned at all.}
{}

% REFUTED
\feedback{The motivation section could be strengthened. While the proposal explains that existing tools are either too technical or not presentation-friendly, it does not fully explain why this problem is important at a broader level. It would help to more clearly identify who specifically benefits from this tool and why current alternatives are insufficient. The proposal currently reads as though the improvement is mostly aesthetic, so emphasizing the educational or practical impact would make the motivation more compelling.}
{
  We believe the motivation section sufficiently explains the importance of the problem and the intended beneficiaries of LLVManim. We feel that the educational and practical impact of LLVManim is clear from the context provided, and, frankly, the improvement is primarily aesthetic.
}

% ADDRESSED - Improved Project Goals section
\feedback{[some objectives] could be more measurable. For example, the proposal mentions performance and trace mapping goals, but it would be helpful to specify what constitutes success more concretely. Adding clearer success criteria would make it easier to evaluate whether the project meets its goals.}
{}

% ADDRESSED - Added IR scope and edge case handling to Ingestion Layer
\feedback{The [architecture] scope could be clarified further. It would be helpful to specify what subset of LLVM IR will be fully supported and how unsupported constructs will be handled. Addressing potential edge cases would make the approach feel more complete and realistic.}
{}

% ADDRESSED - Added deliverables list to project goals
\feedback{The proposal would benefit from a clearer description of final deliverables. While features are described, it is not entirely clear what concrete artifacts will be submitted at the end of the project. Explicitly listing tangible deliverables would strengthen this section.}
{}

% ADDRESSED - Improved timeline detail
\feedback{The timeline is present but somewhat broad in places. Some weekly goals are described generally, which makes it difficult to determine what measurable progress will be achieved each week. Making each milestone more concrete and demonstrable would better align with the rubric's expectations.}
{}

% ADDRESSED - Removed mention of pilot study
\feedback{The proposal mentions evaluation goals such as trace mapping percentages and usability testing, which is good. However, the evaluation methodology could be described in more detail. For example, specifying the number of participants in the pilot study, the types of tasks they will perform, and how results will be measured would make the evaluation plan stronger and more rigorous.}
{}

% ADDRESSED - Added risk evaluation section
\feedback{A clearly labeled risks section would improve the proposal. While potential challenges are implied, explicitly identifying technical and timeline risks and describing mitigation strategies would demonstrate thoughtful project planning.}
{}

\feedbackdivider

% ADDRESSED - Added supported languages
\feedback{I feel [the PythonTutor section] should include the supported languages of PythonTutor, since the usefulness of the tool depends heavily on what languages are supported. It is more useful if it supports languages that are commonly used for education vs if it supports languages that are more obscure.}
{}

% ADDRESSED - Added footnote explaining DOT files
\feedback{What are DOT files? does this mean unix dot files, or does this stand for something?}
{}

% REFUTED
\feedback{What does verbosity mean in this application and how is it different from animation speed. Does the creator of the animation include an audio track, or is it the speed that words appear? If the latter, how and why are words appearing?}
{
  Verbosity in this context refers to the amount of information displayed in the animation, such as whether variable values are shown or only control flow is visualized. It is different from animation speed, which controls how quickly the animation plays. This is a common feature in debugging tools, and we feel it is reasonable to expect readers to understand this distinction.
}

% ADDRESSED - Removed mention of pilot study
\feedback{Go more in depth about what the study entails. Will you be comparing against other similar tools, or just compared to only reading the code? Will it be between subjects or within subjects? Are you only going to collect data about accuracy or will you also collect data about subjective user preference?

Also this study does not seem like it connects well to the stated goal of this project. You are discussing making a tool that will help educational scenarios, but are only evaluating it for use in debugging. I would add a component focusing on user experience generating the images, or evaluating which ones help people to better understand the purpose of code.

There are also no sample graphs or tables for the data that is to be collected.}
{}

% ADDRESSED - Combined core and major features
\feedback{Why separate core features from MVP features when core features are defined as `required to have a usable product at all', which is essentially the MVP?}
{}

% REFUTED
\feedback{How are animations utilized. Is it moving arrows to demonstrate code flow? Will sections of the code be moved on screen only when relevant, and removed afterward. For someone unfamiliar with Manim this is unclear.}
{
  A major portion of this project will be dedicated to exploring precisely this question. While we have some initial mockups and ideas for how the animations will work, these likely will not be finalized until towards the end of the project.
}

% ADDRESSED - Added mention of "Control-Flow Graph" term
\feedback{While I understand that this document is written with the assumption that the reader has a computer science bachelors degree, I think it would help to use the long form of this abbreviation at least once, as was done with LLVM IR.}
{}

% ADDRESSED - Removed mention of "JSON-like" format
\feedback{Why make your own `JSON-like' format instead of just using standard JSON. If it is for data streaming purposes, perhaps consider JSON Lines.}
{}

% ADDRESSED - Added risk evaluation section
\feedback{No discussion of risks or costs.}
{}

% ADDRESSED - Removed mention of pilot study
\feedback{Timeline does not say when the experiment will be performed.}
{}

% ADDRESSED - Added time spent on project
\feedback{Missing time spent on project.}
{}

% ADDRESSED - Added Manim/LLVManim examples
\feedback{The lack of Manim examples or proposed examples for this project makes it difficult for someone unfamiliar with that tool to understand what this tool will produce.}
{}
